% Ch17 WS4 -- dissolution, coupled reactions, and a 10-point FRQ. Block 5.
\documentclass{apchem-worksheet}

\apcunitword{Chapter}
\apcunit{17}
\apcunittitle{Entropy and Free Energy}
\apcdoctitle{Worksheet 4 \textbullet\ Dissolution, Coupling \& FRQ}
\apcsubtitle{Zumdahl \S17.5, 17.10 \textbullet\ free energies add, which is what makes coupling work}

\begin{document}
\wsheader[$\Delta G^\circ = \Delta H^\circ - T\Delta S^\circ$ \textbullet\
$\Delta G^\circ = -RT\ln K$ \textbullet\ $R =
\SI{8.314}{\joule\per\mole\per\kelvin}$. ATP hydrolysis:
$\Delta G^\circ = \SI{-30.5}{\kilo\joule\per\mole}$. $S^\circ$ (J/K$\cdot$mol):
\ce{H2} 131 \textbullet\ \ce{O2} 205 \textbullet\ \ce{H2O(l)} 70.
$\Delta G_f^\circ$ (kJ/mol): \ce{H2O(l)} $-237$.]

\problem[]{Dissolving a salt has three free-energy contributions. Name each
and give the sign of its enthalpy or entropy contribution:}
\begin{parts}
  \item Breaking up the crystal lattice:
        \answerblank[6.4cm]{$\Delta H$ positive, $\Delta S$ positive}
  \item Reorganizing solvent around the ions:
        \answerblank[7.6cm]{$\Delta S$ negative (hydration shells order)}
  \item Attraction between ions and solvent:
        \answerblank[6.4cm]{$\Delta H$ negative}
\end{parts}

\problem[]{The CED says predicting the overall $\Delta G^\circ$ of
dissolution ``can be challenging.'' Explain why, and say what you
\emph{are} expected to do.
\answerlines[4]{The three contributions largely cancel: two large enthalpy
terms of opposite sign, and two entropy terms of opposite sign. The net
free energy is a small difference between big numbers, so its sign cannot
be reliably predicted by inspection. You are expected to identify the
contributions and their individual signs and reason about them, not to
predict solubility from scratch.}}

\problem[]{\ce{NH4NO3} dissolves \emph{endothermically} --- the beaker gets
noticeably cold --- yet it is extremely soluble.}
\begin{parts}
  \item What does the cold beaker tell you about $\Delta H^\circ$?
        \answerblank[3.4cm]{it is positive}
  \item Since the salt dissolves readily, what must be true of
        $\Delta G^\circ$, and therefore of $\Delta S^\circ$?
        \answerlines[3]{$\Delta G^\circ$ must be negative. With
        $\Delta H^\circ$ positive, the only way to achieve that is for
        $T\Delta S^\circ$ to exceed it, so $\Delta S^\circ$ must be
        positive and reasonably large.}
  \item Explain the entropy increase in physical terms.
        \answerlines[2]{A highly ordered ionic crystal breaks up into
        individual ions free to move throughout the solvent, which is a
        large increase in the dispersal of matter.}
  \item What does this example prove about the claim ``exothermic reactions
        are the favored ones''? \answerlines[2]{It disproves it. Enthalpy is
        only one of two terms; here entropy alone drives the process while
        enthalpy opposes it.}
\end{parts}

\problem[]{\ce{NaCl} dissolves with $\Delta H^\circ$ very close to zero, yet
readily. Explain what must be driving it and write the reduced form of the
free-energy expression that applies.
\answerlines[3]{The entropy term is driving it. With
$\Delta H^\circ \approx 0$ the expression reduces to
$\Delta G^\circ \approx -T\Delta S^\circ$, and since dissolution disperses
an ordered lattice into mobile hydrated ions, $\Delta S^\circ > 0$ and
$\Delta G^\circ < 0$.}}

\problem[]{Two ways to drive an unfavorable process.}
\begin{parts}
  \item Name them. \answerlines[2]{Supply energy from an external source, or
        couple the process to a thermodynamically favorable reaction.}
  \item Give the CED's two examples of an external source.
        \answerlines[2]{Electrical energy driving an electrolytic cell or
        charging a battery; light driving the conversion of carbon dioxide
        to glucose in photosynthesis.}
  \item What must two coupled reactions share for the coupling to be real
        rather than bookkeeping?
        \answerblank[5.4cm]{a common intermediate}
\end{parts}

\problem[]{Coupled reaction arithmetic.}
\begin{parts}
  \item Reaction X has $\Delta G^\circ = +9.2$~kJ/mol and is coupled to ATP
        hydrolysis. Find the overall $\Delta G^\circ$ and say whether it is
        favored. \answerblank[5cm]{$-21.3$ kJ/mol; favored}
  \item Reaction Y has $\Delta G^\circ = +50$~kJ/mol. Show that coupling to
        a single ATP hydrolysis is not enough, then find how many are
        needed. \workspace[2.6cm]
        \keyonly{\begin{apcanswer}One ATP: $+50 + (-30.5) = +19.5$~kJ ---
        still unfavorable. Two ATP: $+50 + 2(-30.5) =
        \mathbf{-11.0}~\text{kJ}$ --- favored. So \textbf{two} are
        required.\end{apcanswer}}
  \item Explain what ``ATP is the cell's energy currency'' actually means,
        thermodynamically. \answerlines[3]{Not that ATP stores energy in
        some special bond, but that its hydrolysis has a substantially
        negative $\Delta G^\circ$. Because free energies add, that
        favorable reaction can be coupled to reactions the cell needs which
        are unfavorable on their own, making the sum favorable.}
\end{parts}

\problem[]{Photosynthesis converts \ce{CO2} and water to glucose and oxygen
with $\Delta G^\circ \approx +2870$~kJ per mole of glucose.}
\begin{parts}
  \item Is the reaction thermodynamically favored?
        \answerblank[2.6cm]{no}
  \item Explain how it occurs anyway, and why no law is broken.
        \answerlines[3]{It is driven by an external energy source: absorbed
        sunlight supplies the free energy the reaction lacks. Considering
        the plant together with the light source and surroundings, the
        overall $\Delta G$ is still negative and
        $\Delta S_{\text{univ}} > 0$.}
\end{parts}

\problem[]{\textbf{FRQ (10 points).} Consider the electrolysis of water,
\ce{2H2O(l) -> 2H2(g) + O2(g)}.}
\begin{parts}
  \item Predict the sign of $\Delta S^\circ$ and justify from the equation
        alone. \answerlines[2]{Positive. Three moles of gas are produced
        from a liquid, and gases carry far more entropy than condensed
        phases.}
  \item Calculate $\Delta S^\circ$. \workspace[2.2cm]
        \keyonly{\begin{apcanswer}$[2(131) + 205] - 2(70) = 467 - 140 =
        \mathbf{+327}~\si{\joule\per\kelvin}$\end{apcanswer}}
  \item Calculate $\Delta G^\circ$ at \SI{298}{\kelvin} from
        $\Delta G_f^\circ$ values, and state whether the reaction is
        thermodynamically favored. \workspace[2.4cm]
        \keyonly{\begin{apcanswer}$\Delta G^\circ = 0 - 2(-237) =
        \mathbf{+474}~\text{kJ}$ --- strongly \textbf{not}
        favored\end{apcanswer}}
  \item Given $\Delta H^\circ = +572$~kJ, calculate the temperature above
        which the reaction becomes favored, and comment on whether that is
        a practical route to hydrogen. \workspace[2.8cm]
        \keyonly{\begin{apcanswer}$T = \Delta H^\circ/\Delta S^\circ =
        572/0.327 = \mathbf{1749}~\si{\kelvin}$
        ($\approx\SI{1480}{\celsius}$). Impractical --- maintaining that
        temperature is far harder than supplying the energy
        electrically.\end{apcanswer}}
  \item Explain how hydrogen is obtained from water in practice, in terms
        of this unit's vocabulary. \answerlines[3]{By electrolysis: an
        external electrical energy source drives the thermodynamically
        unfavorable reaction. The electricity supplies the free energy the
        reaction lacks, exactly as light does for photosynthesis.}
\end{parts}

\begin{apcnote}[Rubric --- score yourself, 10 points]
\textbf{(a) 2 pts} 1 positive, 1 justification citing 3 mol of gas formed
from a liquid \textbullet\ \textbf{(b) 2 pts} 1 correct setup with
coefficients, 1 answer $+327$~J/K \textbullet\ \textbf{(c) 2 pts} 1
$\Delta G^\circ = +474$~kJ, 1 states not favored \textbullet\
\textbf{(d) 3 pts} 1 recognizes $T = \Delta H^\circ/\Delta S^\circ$,
1 unit conversion to kJ/K and answer $1749$~K, 1 comment on impracticality
--- an answer that forgets to convert J to kJ earns at most 1 \textbullet\
\textbf{(e) 1 pt} electrolysis identified as an external energy source
driving an unfavorable process.
\end{apcnote}

\begin{spiralstrip}{Chapter 17 \textbullet\ blocks 1--4}
  \item $\Delta G^\circ = -RT\ln K$; if $K > 1$ then $\Delta G^\circ$ is
        \answerblank[2cm]{negative}
  \item Favored but immeasurably slow means:
        \answerblank[3.4cm]{kinetic control}
  \item Crossover temperature $=$ \answerblank[2.6cm]{$\Delta H/\Delta S$}
  \item A catalyst changes the rate but not:
        \answerblank[2cm]{$K$}
  \item $\Delta H^\circ < 0$ with $\Delta S^\circ > 0$ is favored at:
        \answerblank[3.6cm]{all temperatures}
\end{spiralstrip}

\end{document}
